%! Elimination Matrices and Inverse Matrices — Unit 2, Lesson 2.2 — Guided Notes
\documentclass[10pt]{article}
\usepackage{linalg-article}
\usepackage{linalg-boxes}
\newcommand{\TallMath}[1]{$\displaystyle #1\rule[-1.4em]{0pt}{3.2em}$}
\newcommand{\xx}{\mathbf{x}}
\newcommand{\bb}{\mathbf{b}}

\begin{document}

\pageheader{Unit 2, Lesson 2.2}{Guided Notes}

\vspace{0.12in}

\begin{objectivebox}
\textbf{By the end of this lesson, I will be able to\dots}
\begin{itemize}\setlength{\itemsep}{2pt}
  \item write one elimination step as an \emph{elimination matrix} $E$, and check that $EA$ does the step;
  \item \emph{undo} a step with $E^{-1}$ (add the multiplier back), knowing $E^{-1}E=I$;
  \item build a $2\times2$ \emph{inverse} $A^{-1}$ (the $ad-bc$ formula and Gauss--Jordan) and solve with $\xx=A^{-1}\bb$;
  \item recognize a \emph{singular} matrix ($ad-bc=0$) as one with no inverse.
\end{itemize}
\end{objectivebox}

\vspace{0.10in}

\begin{vocabbox}
\small
Fill in each term as we build it together.
\par\vspace{2pt}
\termblanklong{Identity matrix $I$}
\termblanklong{Elimination matrix $E$}
\termblanklong{Inverse matrix $A^{-1}$}
\termblanklong{Gauss--Jordan}
\termblanklong{Singular matrix}
\end{vocabbox}

\begin{hookbox}
A bakery ships the same two gift boxes all season, so its ``ingredients'' matrix never changes ---
only the weekly order $\bb$ does.
\[
  A = \begin{bmatrix} 1 & 3 \\ 2 & 7 \end{bmatrix}, \qquad
  \text{this week } \bb = \begin{bmatrix} 7 \\ 15 \end{bmatrix}.
\]
\begin{itemize}\setlength{\itemsep}{1pt}
  \item In 2.1 we would run elimination from scratch every week. If only $\bb$ changes, what could
        we build \emph{once} to solve every future order in a single step? \writeline
\end{itemize}
\end{hookbox}

\begin{notesbox}{1. One Elimination Step Is a Matrix}
Start with the \textbf{identity matrix} $I=\begin{bsmallmatrix} 1&0\\0&1 \end{bsmallmatrix}$. It leaves
any matrix alone: $IA=A$ and $I\bb=\bb$.

\vspace{3pt}
Now take the bakery system's multiplier, using 2.1's rule: pivot $1$, so $\ell=2\div1=2$. Drop $-\ell$ into the
lower-left spot of $I$ to get the \textbf{elimination matrix}:
\[
  E_{21} = \begin{bmatrix} 1 & 0 \\ -\ell & 1 \end{bmatrix} = \begin{bmatrix} 1 & 0 \\ \blank{0.8cm} & 1 \end{bmatrix}.
\]
Watch $E_{21}A$ do exactly one elimination step --- row~1 untouched, row~2 becomes (row~2)$-2\,$(row~1):
\[
  E_{21}A = \begin{bmatrix} 1 & 0 \\ -2 & 1 \end{bmatrix}\begin{bmatrix} 1 & 3 \\ 2 & 7 \end{bmatrix}
  = \begin{bmatrix} 1 & 3 \\ \blank{0.8cm} & \blank{0.8cm} \end{bmatrix} = U.
\]
The same matrix updates the right-hand side: $E_{21}\bb=\begin{bsmallmatrix} 1&0\\-2&1 \end{bsmallmatrix}\begin{bsmallmatrix} 7\\15 \end{bsmallmatrix}=\begin{bsmallmatrix}\blank{0.8cm}\\[2pt]\blank{0.8cm}\end{bsmallmatrix}$ ---
the triangular system, ready for back-substitution.
\end{notesbox}

\begin{notesbox}{2. Undoing a Step: The Inverse of $E$}
Every step is reversible. $E_{21}$ \emph{subtracted} $2\times$row~1 from row~2; to undo it you
\textbf{add} $2\times$row~1 back. So put $+\ell$ where the $-\ell$ was:
\[
  E_{21}^{-1} = \begin{bmatrix} 1 & 0 \\ +\ell & 1 \end{bmatrix} = \begin{bmatrix} 1 & 0 \\ \blank{0.8cm} & 1 \end{bmatrix}.
\]
Check that it truly undoes $E_{21}$ --- multiplying should return the identity:
\[
  E_{21}^{-1}E_{21} = \begin{bmatrix} 1 & 0 \\ 2 & 1 \end{bmatrix}\begin{bmatrix} 1 & 0 \\ -2 & 1 \end{bmatrix}
  = \begin{bmatrix}\rule{0pt}{1.1em}\blank{0.8cm} & \blank{0.8cm}\\[2pt]\blank{0.8cm} & \blank{0.8cm}\end{bmatrix} = I.
\]
This is what an \textbf{inverse} means: the matrix that multiplies back to $I$.
\end{notesbox}

\begin{notesbox}{3. The Inverse of $A$ --- the Solving Machine}
The \textbf{inverse} $A^{-1}$ is the matrix with $A^{-1}A=I$. It answers the hook: since
$A\xx=\bb$, multiplying both sides by $A^{-1}$ gives
\[
  A^{-1}A\,\xx = A^{-1}\bb \;\Longrightarrow\; \xx = A^{-1}\bb.
\]
Build $A^{-1}$ once, and \emph{every} order is a single multiplication. For a $2\times2$ matrix there
is a formula --- swap the diagonal, negate the off-diagonal, divide by $ad-bc$:
\[
  A = \begin{bmatrix} a & b \\ c & d \end{bmatrix} \;\Longrightarrow\;
  A^{-1} = \frac{1}{ad-bc}\begin{bmatrix} d & -b \\ -c & a \end{bmatrix}.
\]
For $A=\begin{bsmallmatrix} 1&3\\2&7 \end{bsmallmatrix}$: $ad-bc = (1)(7)-(3)(2) = \blank{0.9cm}$, so
\[
  A^{-1} = \frac{1}{1}\begin{bmatrix} 7 & -3 \\ -2 & 1 \end{bmatrix} = \begin{bmatrix} 7 & -3 \\ -2 & 1 \end{bmatrix}.
\]
Now solve this week's order: $\xx = A^{-1}\bb = \begin{bsmallmatrix} 7&-3\\-2&1 \end{bsmallmatrix}\begin{bsmallmatrix} 7\\15 \end{bsmallmatrix} = \begin{bsmallmatrix}\blank{0.8cm}\\[2pt]\blank{0.8cm}\end{bsmallmatrix}$,
so $(x,y)=(\,\blank{1.4cm}\,)$ --- the same answer elimination gives.

\vspace{3pt}
\textbf{Singular warning.} If $ad-bc=0$, the formula divides by zero: there is \textbf{no inverse}.
That is exactly the zero-pivot breakdown from 2.1.
\end{notesbox}

\begin{notesbox}{4. Building $A^{-1}$ by Gauss--Jordan}
The formula is a $2\times2$ shortcut; \textbf{Gauss--Jordan} works for any size. Write $A$ next to
$I$ and row-reduce until the left side becomes $I$ --- the right side turns into $A^{-1}$.
\[
  [\,A \mid I\,] = \left[\begin{array}{cc|cc} 1 & 3 & 1 & 0 \\ 2 & 7 & 0 & 1 \end{array}\right]
  \xrightarrow{\ \text{row}_2 - 2\,\text{row}_1\ }
  \left[\begin{array}{cc|cc} 1 & 3 & 1 & 0 \\ 0 & 1 & \blank{0.7cm} & \blank{0.7cm} \end{array}\right]
\]
Now clear \emph{upward} to finish the left side at $I$:
\[
  \xrightarrow{\ \text{row}_1 - 3\,\text{row}_2\ }
  \left[\begin{array}{cc|cc} 1 & 0 & \blank{0.7cm} & \blank{0.7cm} \\ 0 & 1 & -2 & 1 \end{array}\right]
  = [\,I \mid A^{-1}\,].
\]
Read off $A^{-1}=\begin{bsmallmatrix}\blank{0.7cm} & \blank{0.7cm}\\[2pt]-2 & 1 \end{bsmallmatrix}$ ---
the same inverse as the formula gave.
\end{notesbox}

\boxguard
\begin{practicebox}
Let $A=\begin{bmatrix} 3 & 5 \\ 1 & 2 \end{bmatrix}$ and $\bb=\begin{bmatrix} 11 \\ 4 \end{bmatrix}$.
\begin{enumerate}[leftmargin=1.6em]
  \item Compute $ad-bc$ and write $A^{-1}$ from the formula. \hfill $ad-bc=\blank{1.2cm}$
\begin{work}
  A^{-1} &= \begin{bsmallmatrix} 2&-5\\-1&3 \end{bsmallmatrix}
\end{work}
  \item Use $\xx=A^{-1}\bb$ to solve $A\xx=\bb$. \hfill $x=\blank{1.2cm}$\quad $y=\blank{1.2cm}$
\begin{work}
  \begin{bsmallmatrix} 2&-5\\-1&3 \end{bsmallmatrix}\begin{bsmallmatrix} 11\\4 \end{bsmallmatrix}
      &= \begin{bsmallmatrix} 22-20\\-11+12 \end{bsmallmatrix}
       = \begin{bsmallmatrix} 2\\1 \end{bsmallmatrix}
\end{work}
  \item \textbf{Check} that your $A^{-1}$ is right by multiplying $A^{-1}A$ --- you should get $I$.
\begin{work}
  \begin{bsmallmatrix} 2&-5\\-1&3 \end{bsmallmatrix}\begin{bsmallmatrix} 3&5\\1&2 \end{bsmallmatrix}
      &= \begin{bsmallmatrix} 1&0\\0&1 \end{bsmallmatrix} = I \ \ \text{\checkmark}
\end{work}
\end{enumerate}
\end{practicebox}

\end{document}
